\documentclass[a4paper]{resume} % Use the custom resume.cls style
\usepackage[left=0.4in,top=0.2in,right=0.4in,bottom=0.2in]{geometry} % Document margins

\usepackage{eurosym}
\usepackage{xcolor}

\begin{document}

%-------------------------------------------------------------------------------------------------------------------------------
% EN-TÊTE
%-------------------------------------------------------------------------------------------------------------------------------
\begin{center}
    {\bf \huge Youenn BOGAER} \\
    \vspace{2mm}
    {\bf \large Recherche CDI $-$ Ingénieur IA/ML\\}
    {Disponible à partir de Décembre 2026}
\end{center}

\vspace{3mm}

% -- Colonne gauche --
\begin{minipage}[t]{0.26\textwidth}
\vspace{0pt}
\linespread{0.2}
\small

%----------------------------------------------------------------------------------------
%    Profil SECTION
%----------------------------------------------------------------------------------------
\begin{rLeftColSection}{Profil}
{
    \setlength{\rightskip}{3mm}
    \fontsize{10pt}{25mm}\selectfont
    Futur ingénieur diplômé de CY Tech. Fort d'une formation approfondie en Deep Learning, je me spécialise dans l'entraînement et l'intégration de modèles IA, avec des expériences en LLMs, vision et argentiques. Un socle solide en mathématiques et physique me permet d'aborder des problématiques de recherche appliquées.
    \par
}
\end{rLeftColSection}
\vspace{2mm}
%----------------------------------------------------------------------------------------
%    Contact SECTION
%----------------------------------------------------------------------------------------
\begin{rLeftColSection}{Contact}
\begin{list}{$\cdot$}{\leftmargin=0.5em}
\vspace{1pt}
\itemsep 0.2em
    \item +33(0)6.13.75.00.08
    \item bogaeryouenn@gmail.com
    \item linkedin.com/in/youenn-bogaer-505185221
    \item github.com/YouennBogaer
\end{list}
\end{rLeftColSection}
\vspace{2mm}
%----------------------------------------------------------------------------------------
%    Language and IT skills SECTION
%----------------------------------------------------------------------------------------

\begin{rLeftColSection}{Compétences}
\vspace{2pt}
\textbf{Programmation :}
\begin{list}{}{\leftmargin=0.5em}
\itemsep 0.1em
    \item Python (Avancé) 
    \item SQL, Java, C
\end{list}
\vspace{4pt}
\textbf{IA \& MLOps :}
\begin{list}{}{\leftmargin=0.5em}
\itemsep 0.1em
    \item PyTorch, CrewAI
    \item RAG, Huggingface
    \item Kedro, vLLM
\end{list}
\vspace{4pt}
\textbf{DevOps :}
\begin{list}{}{\leftmargin=0.5em}
\itemsep 0.1em
    \item Kubernetes, Docker
    \item Ansible, Git
\end{list}
\vspace{4pt}
\textbf{Langues :}
\begin{list}{}{\leftmargin=0.5em}
\itemsep 0.1em
    \item Français (Natif)
    \item Anglais (C1$-$TOEIC)
    \item Allemand (B1)
    \item Breton (Natif)
\end{list}
\end{rLeftColSection}

%----------------------------------------------------------------------------------------
%    interests SECTION
%----------------------------------------------------------------------------------------

\end{minipage}
\hfill
\vrule
%----------------------------------------------------------------------------------------
%    Colonne droite
%----------------------------------------------------------------------------------------
\hfill
\begin{minipage}[t]{0.66\textwidth}
\vspace{0pt}
\linespread{0.9}

%----------------------------------------------------------------------------------------
%    Professional experience SECTION
%----------------------------------------------------------------------------------------
\begin{rSection}{Expérience}

\begin{rSubsection}{Stage R\&D Intelligence Artificielle $-$ Armée de Terre, Rennes}{Mai$-$Nov. 2026}{Déploiement d'un SOC assisté par IA}{}
\item Conception et implémentation d'architectures agentiques (CrewAI) pour le triage, l'enrichissement d'alertes et la détection d'attaques.
\item Déploiement multi-modèles \textit{on-premise} et optimisation de l'allocation des ressources matérielles (serveur 90+Go VRAM, vLLM, Kubernetes).
\item Définition des cas d'usage IA, sélection et évaluation des modèles.
\end{rSubsection}

\begin{rSubsection}{Projet de Recherche (PFE) $-$ IRD}{2025$-$2026}{Classification d'images d'herbiers contraint par la morphologie}{}
\item Conception et implémentation sous PyTorch d'une architecture multi-branches intégrant la détection d'organes à la \textit{loss}.
\item Synthèse et rédaction de l'état de l'art scientifique.
\end{rSubsection} 

\begin{rSubsection}{Stage Développeur $-$ Naval Group, Brest}{Été 2024}{Développement sur xWiki et gestion de projet en autonomie}{}
\item Recueil des besoins et développement d'une application de gestion pour les équipes du recrutement (10+ utilisateurs).
\end{rSubsection}

\end{rSection}
%----------------------------------------------------------------------------------------
%    Education SECTION
%----------------------------------------------------------------------------------------

\begin{rSection}{Formation}

\item {\bf CY Tech (ex EISTI), Cergy} \hfill {2023$-$2026}
\item \textit{Grande école d'ingénieurs $-$ Spécialité Intelligence Artificielle} 

\item {\bf Technische Universität Dresden, Allemagne} \hfill {Eté 2025}
\item \textit{Echange Erasmus - Semestre d'été} 

\item {\bf Lycée La Pérouse-Kerichen, Brest} \hfill {2020$-$2023}
\item \textit{CPGE Mathématiques et Physique}

\end{rSection}



%----------------------------------------------------------------------------------------
%    Other projects SECTION
%----------------------------------------------------------------------------------------
\begin{rSection}{Projets académiques}

\item {\bf LLM Multimodal NLP : Image Captioning} \hfill {Décembre 2025}
\item Fine-tuning du modèle LLaVA servi via Ollama. Compréhension fine des architectures hybrides vision/langage.

\item {\bf Computer Vision : Image Segmentation} \hfill {Novembre 2025}
\item Fine-tuning d’un modèle Mask R-CNN pour la segmentation d'instances.

\item {\bf Management $-$ Lauréat Hackathon junior} \hfill {Décembre 2025}
\item Direction et formation d'une équipe de 5 personnes sous forte contrainte de temps.
\end{rSection} 



\end{minipage}

\end{document}